% =====================================================================
% บทที่ 2 ความต้องการของระบบและการติดตั้ง (ผู้อ่าน: ผู้ดูแลระบบ นักพัฒนา)
% =====================================================================
\label{ch:install}

\begin{guide}
ผู้อ่านบทนี้คือผู้ที่ต้องติดตั้งระบบจากซอร์สโค้ดหรือไฟล์ที่ส่ง เช่น ผู้ดูแลระบบ กรรมการ
หรือรุ่นน้องที่รับช่วงต่อ ให้มีเนื้อหาดังนี้
\begin{itemize}
  \item ความต้องการของระบบ แยกฝั่งผู้ใช้งาน (อุปกรณ์ เบราว์เซอร์ ระบบปฏิบัติการ)
        และฝั่งเครื่องแม่ข่าย (หน่วยประมวลผล หน่วยความจำ พื้นที่ดิสก์ ระบบปฏิบัติการ)
        ระบุทั้งค่าขั้นต่ำและค่าที่แนะนำ
  \item ซอฟต์แวร์ที่ต้องติดตั้งก่อน พร้อมรุ่นและคำสั่งตรวจสอบ
  \item ขั้นตอนติดตั้งเรียงตามลำดับ ทุกคำสั่งคัดลอกไปใช้ได้ และบอกผลลัพธ์ที่ควรเห็นหลังขั้นตอนสำคัญ
  \item การตั้งค่าเป็นตาราง (ชื่อ ค่าตัวอย่าง ความหมาย) และค่าที่ต้องเปลี่ยนก่อนใช้งานจริง
  \item วิธีตรวจว่าติดตั้งสำเร็จ และบัญชีเริ่มต้นสำหรับทดสอบ
\end{itemize}
แอปบนมือถือ: ระบุรุ่น Android/iOS วิธีติดตั้งไฟล์ APK หรือจากสโตร์ และสิทธิ์ที่แอปขอ \\
ระบบฮาร์ดแวร์/IoT: ใส่รายการอุปกรณ์พร้อมรุ่น แผนผังการต่อวงจร แหล่งจ่ายไฟ และการอัปโหลดเฟิร์มแวร์ \\
งานวิจัย/โมเดล AI: ระบุรุ่น Python ไลบรารี GPU และตำแหน่งดาวน์โหลดโมเดลหรือชุดข้อมูล
\end{guide}

\section{ความต้องการของระบบ}
\label{sec:requirements}

ระบบเป็นเว็บแอปพลิเคชัน ผู้ใช้งานใช้ผ่านเว็บเบราว์เซอร์โดยไม่ต้องติดตั้งโปรแกรมเพิ่ม
ส่วนเครื่องแม่ข่ายติดตั้งส่วนหลัง ฐานข้อมูล และตัวแปลเอกสาร \LaTeX{} ซึ่งทำงานในคอนเทนเนอร์ Docker

\subsection{ความต้องการฝั่งผู้ใช้งาน}

อุปกรณ์และซอฟต์แวร์ที่ผู้ใช้งานต้องมีแสดงดังตารางที่~\ref{tab:client-req}

\begin{table}[H]
\caption{ความต้องการฝั่งผู้ใช้งาน}
\label{tab:client-req}
\centering
\begin{tabular}{|L{4cm}|L{\dimexpr\linewidth-4cm-4\tabcolsep-3\arrayrulewidth\relax}|}
\hline
\textbf{รายการ} & \textbf{ข้อกำหนด} \\ \hline
เว็บเบราว์เซอร์ & Google Chrome, Microsoft Edge หรือ Mozilla Firefox รุ่นล่าสุด \\ \hline
ความละเอียดหน้าจอ & ตั้งแต่ 1366 $\times$ 768 พิกเซลขึ้นไป (แนะนำ 1920 $\times$ 1080) \\ \hline
บัญชีผู้ใช้ & อีเมลของมหาวิทยาลัย (\cmd{@kku.ac.th}) \\ \hline
การเชื่อมต่อ & อินเทอร์เน็ตตลอดการใช้งาน (การแก้ไขร่วมกันใช้ WebSocket) \\ \hline
\end{tabular}
\end{table}

\subsection{ความต้องการของเครื่องแม่ข่าย}

เครื่องแม่ข่ายสำหรับใช้งานจริงควรมีคุณสมบัติตามตารางที่~\ref{tab:server-req}

\begin{table}[H]
\caption{ความต้องการของเครื่องแม่ข่าย}
\label{tab:server-req}
\centering
\begin{tabular}{|L{4cm}|L{\dimexpr0.5\linewidth-2cm-3\tabcolsep-2\arrayrulewidth\relax}|L{\dimexpr0.5\linewidth-2cm-3\tabcolsep-2\arrayrulewidth\relax}|}
\hline
\textbf{รายการ} & \textbf{ขั้นต่ำ} & \textbf{แนะนำ} \\ \hline
ระบบปฏิบัติการ & Ubuntu 22.04 LTS & Ubuntu 24.04 LTS \\ \hline
หน่วยประมวลผล & 2 vCPU & 4 vCPU \\ \hline
หน่วยความจำ & 4 GB & 8 GB \\ \hline
พื้นที่ดิสก์ (SSD) & 20 GB & 50 GB \\ \hline
เครือข่าย & 100 Mbps & 1 Gbps \\ \hline
\end{tabular}
\end{table}

\begin{note}
อิมเมจ TeX Live ใช้พื้นที่ประมาณ 5 GB และการ Compile แต่ละครั้งใช้หน่วยความจำได้ไม่เกิน 512 MB
จำนวนผู้ใช้ที่ Compile พร้อมกันได้จึงขึ้นกับหน่วยความจำของเครื่องแม่ข่าย
\end{note}

\subsection{ซอฟต์แวร์ที่ต้องติดตั้ง}

ติดตั้งซอฟต์แวร์ในตารางที่~\ref{tab:prereq} บนเครื่องแม่ข่ายก่อนติดตั้งระบบ
ฐานข้อมูล PostgreSQL และ TeX Live ไม่ต้องติดตั้งแยก เพราะทำงานในคอนเทนเนอร์ Docker

\begin{table}[H]
\caption{ซอฟต์แวร์ที่ต้องติดตั้งบนเครื่องแม่ข่าย}
\label{tab:prereq}
\centering
\begin{tabular}{|L{4.2cm}|L{2.6cm}|L{\dimexpr\linewidth-11cm-8\tabcolsep-5\arrayrulewidth\relax}|L{4.2cm}|}
\hline
\textbf{ซอฟต์แวร์} & \textbf{รุ่น} & \textbf{ใช้สำหรับ} & \textbf{คำสั่งตรวจสอบ} \\ \hline
Git & 2.30 ขึ้นไป & ดาวน์โหลดซอร์สโค้ด & \cmd{git --version} \\ \hline
Docker และ Docker Compose & 24 ขึ้นไป & ฐานข้อมูล ส่วนหลัง และ TeX Live & \cmd{docker --version} \\ \hline
Node.js และ npm & 20 LTS / 9 ขึ้นไป & build ส่วนหน้า & \cmd{node --version} \\ \hline
Nginx & 1.18 ขึ้นไป & เว็บเซิร์ฟเวอร์ & \cmd{nginx -v} \\ \hline
\end{tabular}
\end{table}

\section{การติดตั้งระบบ}
\label{sec:install-steps}

หัวข้อนี้ติดตั้งระบบบนเครื่องแม่ข่ายด้วย Docker Compose ซึ่งเริ่มฐานข้อมูลและส่วนหลังพร้อมกัน
ส่วนการติดตั้งบนเครื่องของนักพัฒนาอธิบายในหัวข้อ~\ref{sec:dev-setup}

\begin{steps}
  \item ดาวน์โหลดซอร์สโค้ดไว้ที่โฟลเดอร์ \filename{~/cskkuproject}
\begin{terminal}
cd ~
git clone <repo-url> cskkuproject
cd cskkuproject
\end{terminal}
  \item สร้างไฟล์ตั้งค่าของส่วนหลัง และไฟล์รหัสผ่านสำหรับ Docker Compose
        แล้วแก้ค่าตามหัวข้อ~\ref{sec:env}
\begin{terminal}
cp server/.env.example server/.env
echo "DB_PASSWORD=<รหัสผ่านฐานข้อมูล>" > .env
echo "JWT_SECRET=<คีย์ลับแบบสุ่ม>" >> .env
\end{terminal}
  \item สร้างอิมเมจ TeX Live สำหรับ Compile เอกสาร (ทำครั้งเดียว ใช้เวลานานตามความเร็วอินเทอร์เน็ต)
\begin{terminal}
docker-compose --profile build up texlive
\end{terminal}
  \item เริ่มฐานข้อมูลและส่วนหลัง แล้วตรวจสถานะ
\begin{terminal}
docker-compose up -d
docker-compose ps
\end{terminal}
        \result{บริการ \cmd{db} และ \cmd{api} มีสถานะ \cmd{Up}}
  \item สร้างข้อมูลเริ่มต้นและบัญชีทดสอบ (ครั้งแรกเท่านั้น)
\begin{terminal}
docker-compose exec api npm run db:seed
\end{terminal}
  \item Build ส่วนหน้าแล้วคัดลอกไปยังโฟลเดอร์ของเว็บเซิร์ฟเวอร์
\begin{terminal}
cd client
npm install
npm run build
sudo mkdir -p /var/www/cskkuproject
sudo cp -r dist/. /var/www/cskkuproject/
\end{terminal}
  \item ตั้งค่า Nginx ตามไฟล์ด้านล่าง แล้วโหลดการตั้งค่าใหม่ด้วยคำสั่ง \cmd{sudo nginx -s reload}
        \result{เปิดที่อยู่ของเครื่องแม่ข่ายด้วยเว็บเบราว์เซอร์แล้วพบหน้าเข้าสู่ระบบ}
\end{steps}

\begin{codefile}{/etc/nginx/sites-available/cskkuproject}
server {
    listen 80;
    server_name cskkuproject.example.com;
    root /var/www/cskkuproject;

    location / {
        try_files $uri /index.html;
    }
    location /api/ {
        proxy_pass http://127.0.0.1:3001;
    }
    location /collab {
        proxy_pass http://127.0.0.1:3001;
        proxy_http_version 1.1;
        proxy_set_header Upgrade $http_upgrade;
        proxy_set_header Connection "upgrade";
    }
}
\end{codefile}

\section{การตั้งค่าระบบ}
\label{sec:env}

ไฟล์ \filename{server/.env} กำหนดค่าการทำงานของส่วนหลังตามตารางที่~\ref{tab:env}
ส่วนไฟล์ \filename{.env} ข้างไฟล์ \filename{docker-compose.yml} เก็บรหัสผ่านฐานข้อมูล (\cmd{DB_PASSWORD})
และคีย์ลับ (\cmd{JWT_SECRET}) ที่ Docker Compose ส่งให้ส่วนหลัง

\begin{envtable}[ตัวแปรสภาพแวดล้อมในไฟล์ server/.env][tab:env]
\envvar{DATABASE_URL}{postgres://...:5432/cskkuproject}{ที่อยู่ฐานข้อมูล PostgreSQL ถ้าไม่กำหนด ระบบใช้ SQLite (สำหรับพัฒนา)}
\envvar{JWT_SECRET}{change-this-secret}{คีย์ลงนาม token สุ่มอย่างน้อย 64 ไบต์ ห้ามใช้ค่าตัวอย่างในระบบจริง}
\envvar{JWT_EXPIRES_IN}{7d}{อายุของ token หลังเข้าสู่ระบบ}
\envvar{BCRYPT_ROUNDS}{12}{จำนวนรอบการเข้ารหัสรหัสผ่าน}
\envvar{PORT}{3001}{พอร์ตของส่วนหลัง}
\envvar{NODE_ENV}{production}{\cmd{development} หรือ \cmd{production}}
\envvar{CLIENT_URL}{https://cskkuproject.example.com}{ที่อยู่ของส่วนหน้า (ใช้ตั้งค่า CORS)}
\envvar{DOCKER_IMAGE}{texlive-thai:2023}{อิมเมจที่ใช้ Compile เอกสาร}
\envvar{DOCKER_MEMORY_LIMIT}{536870912}{หน่วยความจำสูงสุดต่อการ Compile (512 MB)}
\envvar{COMPILE_TIMEOUT_MS}{60000}{เวลาสูงสุดต่อการ Compile (มิลลิวินาที)}
\envvar{MAX_FILE_SIZE}{2097152}{ขนาดไฟล์อัปโหลดสูงสุด (2 MB)}
\end{envtable}

\begin{codefile}[language=bash]{server/.env}
# Server
PORT=3001
NODE_ENV=production
CLIENT_URL=https://cskkuproject.example.com

# JWT
JWT_SECRET=change-this-to-a-strong-random-string
JWT_EXPIRES_IN=7d
BCRYPT_ROUNDS=12
\end{codefile}

\begin{tip}
สร้างค่า \cmd{JWT_SECRET} แบบสุ่มได้ด้วยคำสั่งด้านล่าง แล้วคัดลอกผลลัพธ์ไปใส่ในไฟล์ \filename{.env}
\end{tip}
\begin{terminal}
openssl rand -hex 64
\end{terminal}

\section{การตรวจสอบหลังติดตั้ง}
\label{sec:verify}

\begin{steps}
  \item ตรวจว่าส่วนหลังทำงาน
\begin{terminal}
curl http://localhost:3001/api/health
\end{terminal}
        \result{ได้ข้อความที่ขึ้นต้นด้วย \cmd{\{"status":"ok"}}
  \item เปิดระบบด้วยเว็บเบราว์เซอร์ แล้วเข้าสู่ระบบด้วยบัญชีทดสอบในตารางที่~\ref{tab:demo-accounts}
  \item เปิดโครงงานที่อนุมัติแล้ว (สร้างได้ตามหัวข้อ~\ref{sec:proposal}) แล้วคลิก \btn{Compile}
        \result{ไฟล์ PDF แสดงในพื้นที่ตัวอย่าง แสดงว่าอิมเมจ TeX Live พร้อมใช้งาน}
\end{steps}

\begin{table}[H]
\caption{บัญชีสำหรับทดสอบระบบ}
\label{tab:demo-accounts}
\centering
\begin{tabular}{|L{3.4cm}|L{\dimexpr\linewidth-6.8cm-6\tabcolsep-4\arrayrulewidth\relax}|L{3.4cm}|}
\hline
\textbf{บทบาท} & \textbf{อีเมล} & \textbf{รหัสผ่าน} \\ \hline
นักศึกษา & \cmd{somchai.stud@kku.ac.th} & \cmd{password123} \\ \hline
อาจารย์ที่ปรึกษา & \cmd{somchai@kku.ac.th} & \cmd{password123} \\ \hline
\end{tabular}
\end{table}

\begin{warning}
บัญชีในตารางที่~\ref{tab:demo-accounts} สร้างจากคำสั่ง \cmd{npm run db:seed} เพื่อการทดสอบเท่านั้น
ผู้ดูแลระบบต้องลบหรือเปลี่ยนรหัสผ่านก่อนเปิดใช้งานจริง
\end{warning}
